\documentclass[a4paper,fleqn]{cas-sc}
% Elsevier CAS single-column class (cas-sc.cls, cas-common.sty and cas-model2-names.bst,
% version 2.4, are in this folder). The class loads hyperref, booktabs, multirow and array.
\usepackage[authoryear]{natbib}
\usepackage{tabularx}
% allow line breaks at hyphens inside \url (long dataset slugs)
\expandafter\def\expandafter\UrlBreaks\expandafter{\UrlBreaks\do\-}
% Figures are the notebook plots: figures/segmentation/plots/<file>.png when uploaded with
% the project (Overleaf), otherwise read in place from Materials/Notebooks/segmentation/plots/.
\graphicspath{{figures/}{../../Notebooks/}}
% DOIs in the reference list link to https://doi.org (the .bst default is dx.doi.org)
\newcommand*{\doi}[1]{\href{https://doi.org/#1}{\path{#1}}}
\newcolumntype{Y}{>{\raggedright\arraybackslash}X}
\renewcommand{\arraystretch}{1.08}
\emergencystretch=2em
% One plot at [#1 of] the text width, but never taller than #2 of the text height.
\newcommand{\plot}[3][1]{\includegraphics[width=#1\linewidth,height=#2\textheight,keepaspectratio]{#3}}
% Figures and tables are not floats: each is set exactly where it is written, right after
% the paragraph that first cites it, with the cas-sc caption style. A page may end short
% when the next figure or table does not fit on it.
\makeatletter
\ExplSyntaxOn
\NewDocumentEnvironment { herefigure } { }
  {
    \par \addvspace { \intextsep } \noindent
    \begin{minipage}{\linewidth}
    \__reset_fig:
    \dim_set:Nn \l_fig_width_dim { \linewidth }
    \cs_set_eq:NN \@makecaption \__make_fig_caption:nn
    \def \@captype { figure }
    \centering \sffamily \small
  }
  { \end{minipage} \par \addvspace { \intextsep } }
\NewDocumentEnvironment { heretable } { }
  {
    \par \addvspace { \intextsep } \noindent
    \begin{minipage}{\linewidth}
    \__reset_tbl:
    \dim_set:Nn \l_tbl_width_dim { \linewidth }
    \csgdef { tblwidth } { \dim_use:N \l_tbl_width_dim }
    \cs_set_eq:NN \@makecaption \__make_tbl_caption:nn
    \def \@captype { table }
    \centering \sffamily \small
  }
  { \end{minipage} \par \addvspace { \intextsep } }
% cas-common.sty prints an empty "ORCID(s):" footnote when no author has an ORCID iD.
\RenewDocumentCommand \printorcid { }
  {
    \seq_if_empty:NF \g_stm_orcid_seq
      {
        \group_begin:
          \tex_let:D \thefootnote \relax \footnotetext
            { \raggedright \textsc{orcid}(s):\c_space_token \seq_use:Nn \g_stm_orcid_seq { ;~ } }
        \group_end:
      }
  }
\ExplSyntaxOff
\makeatother
\raggedbottom
\hypersetup{
  pdftitle={A Comprehensive Benchmark of Three-Dimensional Segmentation Architectures for Brain Tumour Segmentation},
  pdfauthor={Parth Mital, Ved Lawand, Biji C. L.},
}

\begin{document}
\let\WriteBookmarks\relax

\shorttitle{Benchmark of 3D architectures for brain tumour segmentation}
\shortauthors{P. Mital et al.}

\title[mode=title]{A Comprehensive Benchmark of Three-Dimensional Segmentation Architectures for Brain Tumour Segmentation}

\author[1]{Parth Mital}
\cormark[1]
\ead{parth.mital2023@vitstudent.ac.in}
\author[1]{Ved Lawand}
\ead{ved.lawand2023@vitstudent.ac.in}
\author[1]{Biji C. L.}
\ead{biji.cl@vit.ac.in}

\affiliation[1]{organization={School of Computer Science and Engineering, Vellore Institute of Technology},
  city={Vellore},
  country={India}}

\cortext[1]{Corresponding author}

\begin{abstract}
Brain tumour segmentation from multimodal magnetic resonance imaging (MRI) supports diagnosis, treatment planning and monitoring, but published comparisons of segmentation networks differ in preprocessing, data splits and training budgets. The contribution to artificial intelligence is a controlled benchmark of six three-dimensional architectures, UNet, an nnU-Net-style network, VNet, UNETR, Swin UNETR and TransBTS, implemented from scratch and trained with identical preprocessing, 50 epochs and the same number of optimiser steps on a grade-stratified split of 257 training, 56 validation and 56 test cases from the Brain Tumour Segmentation (BraTS) 2020 dataset. The engineering application is the choice of a network for a deployed segmentation service. Test cases were evaluated at native resolution with eight-flip test-time augmentation, bootstrap confidence intervals and Holm-corrected Wilcoxon signed-rank tests. UNet achieved the highest mean Dice of 0.8199 (95\% confidence interval 0.7826--0.8533), followed by nnU-Net at 0.8192 and VNet at 0.7965; Swin UNETR (0.7462), TransBTS (0.7291) and UNETR (0.7056) trailed, and every pairwise difference was significant. All models segmented the whole tumour well (Dice 0.8714--0.9012), but enhancing-tumour Dice on low-grade gliomas, 27 of 76 of which have no enhancing tumour, fell to between 0.1613 and 0.4237. VNet reached 97\% of the UNet Dice with 1.40 million parameters, 0.53 training hours and 0.0375 seconds per volume, the best Dice per parameter and per training hour in the study. Convolutional networks therefore remain the stronger choice under this data size and training budget.
\end{abstract}

\begin{keywords}
Brain tumour segmentation \sep 3D UNet \sep Transformer \sep BraTS \sep Medical image analysis \sep Deep learning
\end{keywords}

\maketitle

\section{Introduction}
\label{sec:intro}

Gliomas are the most common and aggressive primary brain tumours, accounting for approximately 80\% of all malignant brain tumours worldwide \citep{bakas2017}. These tumours arise from glial cells and are classified into low-grade gliomas (LGG) and high-grade gliomas (HGG), with the latter being more invasive and associated with poorer patient outcomes. Accurate segmentation of these tumours from magnetic resonance imaging (MRI) scans is essential for several clinical tasks including diagnosis and grading, treatment planning, disease monitoring, and prognosis prediction. The BraTS challenge, established in 2012, has become the premier benchmark for evaluating brain tumour segmentation algorithms. It provides multimodal MRI scans consisting of four complementary sequences: T1-weighted (T1), T1-contrast-enhanced (T1CE), T2-weighted (T2), and fluid-attenuated inversion recovery (FLAIR). Alongside these scans, the dataset includes expert-annotated segmentation masks delineating three tumour sub-regions: whole tumour (WT), tumour core (TC), and enhancing tumour (ET) \citep{menze2015}. Manual delineation of brain tumours is time-consuming, requires specialised expertise, and suffers from inter-rater variability. Deep learning based segmentation methods have demonstrated the potential to automate this process, thereby reducing radiologist workload and providing consistent, reproducible results. However, the rapid proliferation of architectural innovations from convolutional neural networks to transformer based models has created a fragmented landscape where practitioners struggle to identify the most suitable approach for their specific clinical or research settings.

Despite significant advances in deep learning for medical image segmentation, several persistent challenges remain unaddressed in the existing literature. First, the proliferation of architectures from classic UNets to complex transformer hybrids creates confusion about which approach is optimal for clinical deployment. Researchers often introduce novel architectures without rigorous comparison against established baselines under controlled conditions. Second, many published results use different preprocessing pipelines, training protocols, data splits, and evaluation metrics, making fair comparisons across studies difficult or impossible. Third, transformer based models promise better global context modelling through self-attention mechanisms, but they require substantial computational resources, larger datasets, and careful hyperparameter tuning, which raises questions about their practical utility in resource-constrained clinical environments. Fourth, medical datasets are typically limited in size due to privacy concerns, annotation costs, and disease rarity, so architectures with tens or hundreds of millions of parameters are prone to overfitting. Fifth, most existing comparisons focus on mean Dice scores while ignoring boundary accuracy, statistical significance, behaviour on clinically distinct subgroups, uncertainty, computational efficiency and memory footprint.

This study addresses these gaps with the following objectives. First, to establish a fair comparison by implementing six representative architectures, namely UNet, an nnU-Net-style network, VNet, UNETR, Swin UNETR, and TransBTS, from scratch in PyTorch and training them with one code path, one split, one effective batch size and one training length. Second, to evaluate every model on the original image grid of each test case with per-case Dice, boundary distance, sensitivity and specificity, and to quantify differences with bootstrap confidence intervals and paired significance tests. Third, to report results separately for high-grade and low-grade gliomas. Fourth, to analyse the trade-offs between segmentation accuracy, parameter count, training time, inference latency and memory, together with ensembles and test-time-augmentation uncertainty. Fifth, to release the notebook, the split and the per-case outputs so that the comparison can be reproduced and extended. The remainder of this paper is organised as follows. Section~\ref{sec:lit} reviews the evolution of segmentation architectures. Section~\ref{sec:methods} details the dataset, preprocessing, model implementations, training configuration and evaluation protocol. Section~\ref{sec:results} presents and discusses the results. Section~\ref{sec:limitations} discusses limitations and future work, Section~\ref{sec:ethics} ethical implications, and Section~\ref{sec:conclusion} concludes.

\section{Literature Review}
\label{sec:lit}

The evolution of medical image segmentation has been marked by several foundational architectures. The UNet architecture, introduced by \citet{ronneberger2015}, revolutionised medical image segmentation by introducing a symmetric encoder-decoder structure with skip connections. The encoder progressively downsamples the input to capture context, while the decoder upsamples to recover spatial details. The 3D UNet, proposed by \citet{cicek2016}, extended this architecture to volumetric data by replacing 2D convolutions with 3D counterparts, enabling direct processing of 3D MRI volumes. The VNet, developed by \citet{milletari2016}, introduced residual connections allowing deeper networks without gradient degradation, and PReLU activations that improve gradient flow. The nnU-Net, presented by \citet{isensee2021}, represented a paradigm shift by automating architecture design based on dataset characteristics. Rather than a fixed architecture, nnU-Net analyses the dataset fingerprint, including image sizes, voxel spacings, and class imbalances, to configure network depth and width, patch size, batch size, and training schedule.

Parallel to CNN innovations, transformer-based architectures have emerged. Vision Transformers (ViT), introduced by \citet{dosovitskiy2021}, demonstrated that pure transformer architectures could match or exceed CNN performance on image classification tasks by modelling long-range dependencies through self-attention mechanisms. TransUNet \citep{chen2021} placed a transformer encoder on top of a CNN feature extractor for 2D medical segmentation. UNETR, proposed by \citet{hatamizadeh2022unetr}, integrated a ViT encoder with a CNN decoder for 3D segmentation. Swin UNETR, developed by \citet{hatamizadeh2022swin}, replaced the standard ViT with a Swin Transformer backbone, using window-based self-attention and shifted windows to reduce computational complexity. TransBTS, introduced by \citet{wang2021}, took a hybrid approach, placing transformer layers only at the bottleneck of a CNN encoder-decoder. Transformers are known to need large datasets or distillation to train well from scratch \citep{touvron2021}, and several medical variants modify attention to cope with limited data \citep{valanarasu2021,cao2023}.

Despite the proliferation of architectures, several gaps remain: lack of systematic comparison under controlled conditions; implementation variability in preprocessing, data splits and training protocols; inadequate reporting of statistical significance, boundary accuracy, subgroup performance and computational cost; and insufficient code and documentation for reproduction. This study addresses these gaps by implementing all six architectures from scratch in one notebook with identical preprocessing and training budget, evaluating them per case at native resolution with confidence intervals and significance tests, reporting results by tumour grade, measuring latency and memory, and releasing the notebook and its outputs.

\section{Methods and Materials}
\label{sec:methods}

\subsection{Dataset}
\label{sec:dataset}

The BraTS 2020 training dataset \citep{menze2015,bakas2017,bakas2018} was used, in the Kaggle distribution \url{https://www.kaggle.com/datasets/awsaf49/brats20-dataset-training-validation}. Segmentation masks were located by name pattern, which also finds the misnamed mask of case 355, so all 369 training cases have the four sequences and a mask: 293 HGG and 76 LGG cases, with the grade taken from the dataset's \texttt{name\_mapping.csv}. Each case contains four MRI sequences. The T1-weighted sequence provides standard anatomical imaging with good grey-white matter contrast. The T1 contrast-enhanced sequence uses a gadolinium-based contrast agent to highlight areas of blood-brain barrier breakdown, which typically correspond to active tumour regions. The T2-weighted sequence visualises oedema and tumour infiltration into surrounding brain parenchyma. The FLAIR sequence suppresses cerebrospinal fluid signals, making peritumoural oedema and non-enhancing tumour components more conspicuous. Ground truth masks contain label 1 for necrotic and non-enhancing tumour core, label 2 for peritumoural oedema, and label 4 for enhancing tumour. They are converted into three nested binary regions: whole tumour from all non-zero labels, tumour core from labels 1 and 4, and enhancing tumour from label 4. One case with its sequences and regions is shown in Figure~\ref{fig:modalities}.

\begin{herefigure}
\centering
\plot{0.4}{segmentation/plots/02_modalities_and_masks.png}
\caption{One BraTS 2020 case: the four MRI sequences and the three nested tumour regions.}
\label{fig:modalities}
\end{herefigure}

The two grades differ in a way that matters for evaluation (Table~\ref{tab:volumes}). LGG whole tumours are as large as HGG ones, with a median of 86.32 against 91.13~cm$^3$, but their enhancing tumour is much smaller, with a median of 0.82 against 19.46~cm$^3$, and 27 of the 76 LGG cases (35.5\%) have no enhancing tumour at all, whereas every HGG case has some.

\begin{heretable}
\centering
\caption{Reference region volumes in cm$^3$ at native resolution, by tumour grade (293 HGG, 76 LGG cases).}
\label{tab:volumes}
\small
\begin{tabular*}{\textwidth}{@{\extracolsep{\fill}}llccccc@{}}
\toprule
Region & Grade & Mean & SD & Min & Median & Max \\
\midrule
\multirow{2}{*}{WT} & HGG & 96.44 & 53.78 & 7.28 & 91.13 & 256.88 \\
 & LGG & 111.54 & 76.72 & 17.08 & 86.32 & 361.78 \\
\multirow{2}{*}{TC} & HGG & 37.66 & 29.26 & 0.56 & 32.54 & 126.72 \\
 & LGG & 57.85 & 49.66 & 0.83 & 44.59 & 190.19 \\
\multirow{2}{*}{ET} & HGG & 23.39 & 18.50 & 0.11 & 19.46 & 111.25 \\
 & LGG & 5.46 & 13.53 & 0.00 & 0.82 & 89.55 \\
\bottomrule
\end{tabular*}
\end{heretable}

The distributions are plotted in Figure~\ref{fig:volumes}.

\begin{herefigure}
\centering
\plot{0.3}{segmentation/plots/01_region_volumes.png}
\caption{Reference WT, TC and ET volumes by tumour grade.}
\label{fig:volumes}
\end{herefigure}

\subsection{Preprocessing and Split}
\label{sec:preproc}

Each sequence was normalised by a per-volume z-score computed over its non-zero (foreground) voxels only, so that background does not distort the intensity distribution. Images were resampled from their native $240\times 240\times 155$ grid to $128\times 128\times 128$ voxels with trilinear interpolation, and masks with nearest-neighbour interpolation, to fit the networks within the memory of one GPU. The input stacks FLAIR, T1, T1CE and T2 as four channels, in the order used by the web application that deploys the VNet model. The preprocessed volumes were cached once.

A fixed split with seed 42, stratified by tumour grade, divided the cases into 257 training, 56 validation and 56 test cases (Table~\ref{tab:splits}), and every model used the same split.

\begin{heretable}
\centering
\caption{Grade-stratified split of the 369 cases.}
\label{tab:splits}
\small
\begin{tabular*}{\textwidth}{@{\extracolsep{\fill}}lccc@{}}
\toprule
Split & HGG & LGG & Total \\
\midrule
Training & 204 & 53 & 257 \\
Validation & 45 & 11 & 56 \\
Test & 44 & 12 & 56 \\
\textbf{Total} & \textbf{293} & \textbf{76} & \textbf{369} \\
\bottomrule
\end{tabular*}
\end{heretable}

Training augmentation, applied on the GPU, consisted of random flips along each axis and a per-channel intensity scale between 0.9 and 1.1 and shift between $-0.1$ and $0.1$ (Figure~\ref{fig:augmentation}).

\begin{herefigure}
\centering
\plot{0.4}{segmentation/plots/03_augmentation.png}
\caption{Training augmentation applied to one training case: random flips along each axis and per-channel intensity scale and shift.}
\label{fig:augmentation}
\end{herefigure}

\subsection{Model Architectures}
\label{sec:models}

All six networks take the four-channel $128^3$ volume and output three logits per voxel, one per region; every model returns its main output first, followed by any deep-supervision outputs. Parameter counts and estimated GFLOPs per volume were measured by the notebook and are compared with the measured cost of every model in Section~\ref{sec:efficiency}.

\textbf{UNet.} A 3D UNet with four resolution levels of 32, 64, 128 and 256 channels and three max-pooling steps. Each level has two $3\times 3\times 3$ convolutions with affine instance normalisation and LeakyReLU (slope 0.01); the decoder uses transposed convolutions and concatenated skip connections, with deep-supervision outputs at decoder stages 2 and 3. It has 5.60 million parameters.

\textbf{nnU-Net-style network.} A fixed full-resolution architecture in the style of nnU-Net, not the self-configuring nnU-Net pipeline: five stages of 32, 64, 128, 256 and 320 channels with stride-2 convolutions for downsampling, the same convolution blocks as UNet, and deep supervision at four decoder stages. It has 16.55 million parameters.

\textbf{VNet.} A residual VNet with stages of 16, 32, 64 and 128 channels. Residual blocks of one or two $3\times 3\times 3$ convolutions with affine instance normalisation and PReLU use identity shortcuts or $1\times 1\times 1$ projections; downsampling uses stride-2 convolutions and upsampling transposed convolutions, with deep-supervision outputs at two decoder stages. Its module names match the network class of the web application, so the trained weights load there unchanged. It has 1.40 million parameters.

\textbf{UNETR.} A ViT-Base encoder on $16\times 16\times 16$ patches, giving 512 tokens of dimension 768, with 12 pre-norm transformer layers of 12 heads, an MLP dimension of 3,072 and dropout 0.1. Token maps from layers 3, 6, 9 and 12 feed a convolutional decoder with transposed convolutions and blocks of instance normalisation and ReLU, plus a full-resolution convolutional skip from the input. It has 137.60 million parameters, the largest in the study.

\textbf{Swin UNETR.} A Swin Transformer encoder on $2\times 2\times 2$ patches with an embedding of 48 channels doubled at each of four stages (48, 96, 192, 384), $7\times 7\times 7$ windows, 3, 6, 12 and 24 heads, two blocks per stage alternating regular and shifted windows, and a learnable relative position bias per head. The convolutional decoder uses skips from every stage and a full-resolution convolutional skip from the input. It has 10.52 million parameters.

\textbf{TransBTS.} A hybrid with a residual convolutional encoder of 32, 64, 128, 256 and 512 channels and four stride-2 downsamplings, four transformer layers of 8 heads and dimension 512 at the bottleneck with a learnable position embedding for the $8\times 8\times 8$ tokens, and a residual decoder with skip connections and deep-supervision outputs at two decoder stages. It has 35.81 million parameters.

The set covers the standard encoder-decoder baseline (UNet), a deeper nnU-Net-style variant, a compact residual design (VNet), a pure transformer encoder (UNETR), a hierarchical windowed transformer (Swin UNETR) and a hybrid with transformer layers only at the bottleneck (TransBTS).

\subsection{Training Configuration}
\label{sec:training}

The loss for each region was the sum of a soft Dice loss computed over the batch and binary cross-entropy. Deep-supervision outputs were compared with masks downsampled by nearest-neighbour interpolation and weighted 1.0, 0.5 and 0.25 for UNet, VNet and TransBTS, and 1.0, 0.5, 0.25 and 0.125 for the nnU-Net-style network. Optimisers followed the architecture family (Table~\ref{tab:hyperparameters}). The three convolutional networks used SGD with Nesterov momentum 0.99, a learning rate of 0.01, weight decay $3\times 10^{-5}$ and polynomial decay with power 0.9 applied per optimiser step. UNETR, Swin UNETR and TransBTS used AdamW with a learning rate of $10^{-4}$, weight decay $10^{-5}$, a linear warm-up over three epochs and cosine decay to 1\% of the initial rate, also per step.

\begin{heretable}
\centering
\caption{Training settings by model. All models used 50 epochs, an effective batch of 4 volumes (3,250 optimiser steps), gradient clipping at 1.0 and half-precision autocasting. SGD uses Nesterov momentum 0.99 and polynomial decay with power 0.9; AdamW uses cosine decay after a 3-epoch warm-up. DS: deep-supervision weights; GC: gradient checkpointing.}
\label{tab:hyperparameters}
\small
\begin{tabular*}{\textwidth}{@{\extracolsep{\fill}}lccccccc@{}}
\toprule
Model & Optimiser & LR & Weight decay & Schedule & Micro-batch & GC & DS \\
\midrule
UNet & SGD & 0.01 & $3\times 10^{-5}$ & Polynomial & 2 & No & 1, 0.5, 0.25 \\
nnU-Net & SGD & 0.01 & $3\times 10^{-5}$ & Polynomial & 2 & No & 1, 0.5, 0.25, 0.125 \\
VNet & SGD & 0.01 & $3\times 10^{-5}$ & Polynomial & 2 & No & 1, 0.5, 0.25 \\
UNETR & AdamW & $10^{-4}$ & $10^{-5}$ & Warm-up, cosine & 2 & Yes & 1 \\
Swin UNETR & AdamW & $10^{-4}$ & $10^{-5}$ & Warm-up, cosine & 1 & Yes & 1 \\
TransBTS & AdamW & $10^{-4}$ & $10^{-5}$ & Warm-up, cosine & 2 & No & 1, 0.5, 0.25 \\
\bottomrule
\end{tabular*}
\end{heretable}

Every model trained for 50 epochs with an effective batch of 4 volumes, obtained with micro-batches of 2 accumulated over two steps, or micro-batches of 1 accumulated over four steps for Swin UNETR. All models therefore took the same 3,250 optimiser steps. Training used half-precision autocasting with gradient scaling and gradient clipping at a norm of 1.0; UNETR and Swin UNETR used gradient checkpointing to fit in memory. After every epoch the validation cases were evaluated, and the checkpoint with the best validation mean Dice was kept; no early stopping was applied. The random seed was 42 and deterministic kernels were not enforced, so re-runs should reproduce the results closely but not bit for bit.

Two worker processes, one per GPU, trained different models in parallel on Kaggle with two NVIDIA Tesla T4 GPUs, four CPU cores and 33.7~GB of system memory, using Python 3.13.15, PyTorch 2.11.0 and CUDA 12.8. All six runs finished, in 515.0 minutes of training.

\subsection{Evaluation}
\label{sec:evaluation}

\textbf{Post-processing.} After training, the threshold for removing small enhancing-tumour predictions was tuned on the validation cases for each model, choosing among 0, 50, 100, 200, 300, 500 and 800 voxels at native resolution. At test time, sigmoid probabilities of the main output were averaged over the eight combinations of axis flips (test-time augmentation, TTA), resampled to the native grid with trilinear interpolation and thresholded at 0.5. The region hierarchy was enforced, so that ET lies inside TC and TC inside WT, and an ET prediction smaller than the tuned threshold was removed.

\textbf{Metrics.} For every test case and region, the notebook computes the Dice similarity coefficient, scored as 1 when both prediction and reference are empty; the 95th-percentile symmetric Hausdorff distance (HD95) in millimetres using the voxel spacing, set to 0 when both are empty and to the BraTS value of 373.13~mm when only one is empty; and voxel-wise sensitivity and specificity. The mean Dice of a case is the average over the three regions. Model-level Dice is the mean over cases with a percentile bootstrap 95\% confidence interval from 2,000 resamples; HD95 is summarised by its median because empty-region cases contribute the fixed value of 373.13~mm.

\textbf{Comparisons.} Models were compared with Wilcoxon signed-rank tests on per-case mean Dice, paired by case, with Holm correction \citep{holm1979}. Results were also computed separately for HGG and LGG test cases. Two ensembles were evaluated in the same way: all six models, and the three convolutional networks, each averaging member probabilities and using the median of the members' ET thresholds. The voxel-wise standard deviation across the eight TTA flips, averaged inside the predicted whole tumour, served as a per-case uncertainty and was correlated with per-case mean Dice by Spearman's rank correlation.

\textbf{Efficiency.} Latency per volume was measured at batch size 1 under half-precision autocasting, with and without TTA, as the mean and standard deviation over timed passes after warm-up; peak training and inference memory were recorded on the GPU.

\section{Results and Discussion}
\label{sec:results}

\subsection{Training and Model Selection}
\label{sec:selection}

Table~\ref{tab:selection} lists the epoch at which each model reached its best validation mean Dice and the tuned ET threshold. The convolutional networks reached the highest validation Dice, 0.7845 for UNet and 0.7824 for the nnU-Net-style network, whose best epoch was the last, 50. The transformer models peaked earlier, Swin UNETR at epoch 33 and UNETR at 39, at lower validation Dice. Validation tuning removed small ET predictions below 100 to 800 voxels for four models and chose no filtering for Swin UNETR and TransBTS.

\begin{heretable}
\centering
\caption{Model selection on the 56 validation cases.}
\label{tab:selection}
\small
\begin{tabular*}{\textwidth}{@{\extracolsep{\fill}}lccc@{}}
\toprule
Model & Best epoch & Best validation Dice & ET threshold (voxels) \\
\midrule
UNet & 45 & 0.7845 & 100 \\
nnU-Net & 50 & 0.7824 & 200 \\
VNet & 44 & 0.7598 & 200 \\
Swin UNETR & 33 & 0.7235 & 0 \\
TransBTS & 45 & 0.7102 & 0 \\
UNETR & 39 & 0.6960 & 800 \\
\bottomrule
\end{tabular*}
\end{heretable}

The training and validation curves of every model are shown in Figure~\ref{fig:curves}.

\begin{herefigure}
\centering
\plot{0.62}{segmentation/plots/04_training_curves.png}
\caption{Training and validation loss and validation Dice per region for every model.}
\label{fig:curves}
\end{herefigure}

\subsection{Test Results}
\label{sec:test}

Table~\ref{tab:seg-performance} reports the test results at native resolution. UNet achieved the highest mean Dice of 0.8199, followed closely by the nnU-Net-style network at 0.8192 and then VNet at 0.7965. The transformer-based models followed at a clear distance: Swin UNETR at 0.7462, the hybrid TransBTS at 0.7291 and UNETR at 0.7056. The bootstrap confidence intervals of the two leading CNNs, 0.7826--0.8533 and 0.7837--0.8511, overlap almost completely, whereas the intervals of all three transformer-based models lie entirely below them.

\begin{heretable}
\centering
\caption{Test results on the 56 test cases at native resolution, with eight-flip TTA and tuned ET post-processing. Dice is the mean over cases, with a bootstrap 95\% confidence interval for the mean over regions; HD95 is the median over cases in millimetres; sensitivity is the mean over cases.}
\label{tab:seg-performance}
\scriptsize
\setlength{\tabcolsep}{3pt}
\begin{tabular*}{\textwidth}{@{\extracolsep{\fill}}lccccccccccc@{}}
\toprule
 & & \multicolumn{3}{c}{Dice} & \multicolumn{3}{c}{HD95 (mm)} & \multicolumn{3}{c}{Sensitivity} \\
\cmidrule(lr){3-5}\cmidrule(lr){6-8}\cmidrule(l){9-11}
Model & Mean Dice [95\% CI] & WT & TC & ET & WT & TC & ET & WT & TC & ET \\
\midrule
UNet & 0.8199 [0.7826, 0.8533] & 0.9012 & 0.8300 & 0.7286 & 4.24 & 3.24 & 3.67 & 0.9059 & 0.8797 & 0.8486 \\
nnU-Net & 0.8192 [0.7837, 0.8511] & 0.9002 & 0.8309 & 0.7264 & 4.24 & 3.46 & 3.80 & 0.9070 & 0.8831 & 0.8538 \\
VNet & 0.7965 [0.7497, 0.8340] & 0.8865 & 0.7972 & 0.7058 & 4.06 & 4.18 & 3.46 & 0.8899 & 0.8697 & 0.8278 \\
Swin UNETR & 0.7462 [0.7134, 0.7772] & 0.8714 & 0.7495 & 0.6176 & 5.10 & 6.70 & 6.28 & 0.9619 & 0.9316 & 0.9407 \\
TransBTS & 0.7291 [0.6876, 0.7661] & 0.8743 & 0.7488 & 0.5642 & 4.42 & 6.58 & 5.96 & 0.9468 & 0.9384 & 0.9475 \\
UNETR & 0.7056 [0.6700, 0.7384] & 0.8758 & 0.6692 & 0.5720 & 5.05 & 11.29 & 9.15 & 0.9492 & 0.9481 & 0.9520 \\
\bottomrule
\end{tabular*}
\end{heretable}

The ranking is decided by the tumour core and the enhancing tumour. Whole-tumour Dice varies little, from 0.8714 to 0.9012, and every model's median WT boundary error is about 4 to 5~mm. TC Dice ranges from 0.6692 to 0.8309 and ET Dice from 0.5642 to 0.7286, and the median HD95 of UNETR on TC reaches 11.29~mm. The transformer-based models are more sensitive than the CNNs on every region, with ET sensitivity of 0.9407 to 0.9520 against 0.8278 to 0.8538, while their lower Dice shows that this comes from over-segmentation; voxel-wise specificity is above 0.997 for every model and region. Per-case distributions are shown in Figure~\ref{fig:distributions}.

\begin{herefigure}
\centering
\plot{0.4}{segmentation/plots/07_per_case_distributions.png}
\caption{Per-case Dice and HD95 for every region and model, including the two ensembles.}
\label{fig:distributions}
\end{herefigure}

Representative predictions of every model are shown in Figure~\ref{fig:qualitative}.

\begin{herefigure}
\centering
\plot{0.3}{segmentation/plots/08_qualitative_comparison.png}
\caption{The test cases with the median and the lowest all-model-ensemble Dice: reference and every model's prediction on FLAIR at the axial slice with the largest reference tumour (oedema-only WT green, TC yellow, ET red).}
\label{fig:qualitative}
\end{herefigure}

All differences between models are significant after Holm correction (Table~\ref{tab:wilcoxon}, Figure~\ref{fig:wilcoxon}). The smallest is between UNet and the nnU-Net-style network, a median per-case difference of 0.0023 Dice with an adjusted $p$ of 0.042, which is statistically detectable but practically negligible; the same holds for Swin UNETR against TransBTS (0.0064, $p=0.042$). UNet leads VNet by a median of 0.0119 and every transformer-based model by 0.0738 to 0.1139.

\begin{heretable}
\centering
\caption{Pairwise Wilcoxon signed-rank tests on per-case mean Dice, Holm-corrected. A positive median difference favours the first model. Values below $10^{-5}$ are shown as $<10^{-5}$.}
\label{tab:wilcoxon}
\small
\begin{tabular*}{\textwidth}{@{\extracolsep{\fill}}llcc@{}}
\toprule
Model A & Model B & Median difference & Holm $p$ \\
\midrule
UNet & nnU-Net & 0.0023 & 0.0420 \\
UNet & VNet & 0.0119 & $<10^{-5}$ \\
UNet & UNETR & 0.1139 & $<10^{-5}$ \\
UNet & Swin UNETR & 0.0738 & $<10^{-5}$ \\
UNet & TransBTS & 0.0789 & $<10^{-5}$ \\
nnU-Net & VNet & 0.0100 & $1\times 10^{-5}$ \\
nnU-Net & UNETR & 0.1066 & $<10^{-5}$ \\
nnU-Net & Swin UNETR & 0.0687 & $<10^{-5}$ \\
nnU-Net & TransBTS & 0.0795 & $<10^{-5}$ \\
VNet & UNETR & 0.0858 & $<10^{-5}$ \\
VNet & Swin UNETR & 0.0604 & $<10^{-5}$ \\
VNet & TransBTS & 0.0666 & $<10^{-5}$ \\
UNETR & Swin UNETR & $-0.0359$ & $<10^{-5}$ \\
UNETR & TransBTS & $-0.0194$ & 0.0002 \\
Swin UNETR & TransBTS & 0.0064 & 0.0420 \\
\bottomrule
\end{tabular*}
\end{heretable}

\begin{herefigure}
\centering
\plot[0.6]{0.4}{segmentation/plots/05_wilcoxon.png}
\caption{Holm-corrected Wilcoxon signed-rank p-values between models.}
\label{fig:wilcoxon}
\end{herefigure}

\subsection{Results by Tumour Grade}
\label{sec:grade}

Table~\ref{tab:grade} separates the 44 HGG and 12 LGG test cases. On HGG, UNet and the nnU-Net-style network reach mean Dice of 0.8629 and 0.8632, with ET Dice of about 0.81. On LGG, every model drops sharply: mean Dice falls to between 0.5616 and 0.6624, and ET Dice to between 0.1613 (TransBTS) and 0.4237 (UNet). WT Dice, in contrast, is slightly higher on LGG than on HGG for every model.

\begin{heretable}
\centering
\caption{Test Dice by tumour grade (44 HGG, 12 LGG cases), with the bootstrap 95\% confidence interval of the mean Dice.}
\label{tab:grade}
\small
\begin{tabular*}{\textwidth}{@{\extracolsep{\fill}}llccccc@{}}
\toprule
Model & Grade & WT & TC & ET & Mean Dice & 95\% CI \\
\midrule
\multirow{2}{*}{UNet} & HGG & 0.8984 & 0.8786 & 0.8117 & 0.8629 & [0.8300, 0.8889] \\
 & LGG & 0.9115 & 0.6520 & 0.4237 & 0.6624 & [0.6055, 0.7253] \\
\multirow{2}{*}{nnU-Net} & HGG & 0.8985 & 0.8785 & 0.8124 & 0.8632 & [0.8357, 0.8865] \\
 & LGG & 0.9065 & 0.6561 & 0.4112 & 0.6579 & [0.6032, 0.7175] \\
\multirow{2}{*}{VNet} & HGG & 0.8819 & 0.8412 & 0.7923 & 0.8385 & [0.7965, 0.8716] \\
 & LGG & 0.9032 & 0.6359 & 0.3885 & 0.6425 & [0.5750, 0.7131] \\
\multirow{2}{*}{Swin UNETR} & HGG & 0.8620 & 0.7859 & 0.7262 & 0.7913 & [0.7666, 0.8147] \\
 & LGG & 0.9061 & 0.6161 & 0.2193 & 0.5805 & [0.5249, 0.6299] \\
\multirow{2}{*}{TransBTS} & HGG & 0.8665 & 0.7837 & 0.6740 & 0.7748 & [0.7389, 0.8073] \\
 & LGG & 0.9029 & 0.6206 & 0.1613 & 0.5616 & [0.5043, 0.6120] \\
\multirow{2}{*}{UNETR} & HGG & 0.8686 & 0.6843 & 0.6773 & 0.7434 & [0.7100, 0.7736] \\
 & LGG & 0.9021 & 0.6137 & 0.1857 & 0.5672 & [0.5037, 0.6325] \\
\bottomrule
\end{tabular*}
\end{heretable}

Figure~\ref{fig:grade} plots the per-region Dice of both grades. The pattern is consistent with the data: many LGG cases have little or no enhancing tumour (Table~\ref{tab:volumes}), and any false-positive ET voxel on a case without enhancing tumour gives that case an ET Dice of zero. The ET post-processing tuned on the validation set reduces but does not remove this effect, and the two models that chose no ET filtering, Swin UNETR and TransBTS, have two of the three lowest LGG ET scores. The CNN advantage is present in both grades.

\begin{herefigure}
\centering
\plot{0.3}{segmentation/plots/06_dice_by_grade.png}
\caption{Test Dice per region for HGG and LGG cases.}
\label{fig:grade}
\end{herefigure}

\subsection{Ensembles and Uncertainty}
\label{sec:ensembles}

An ensemble of the three CNNs reached a mean Dice of 0.8178 (95\% CI 0.7782--0.8522; WT 0.9010, TC 0.8262, ET 0.7262), and an ensemble of all six models 0.7991 (0.7631--0.8309). Neither exceeded UNet alone, so averaging with weaker models diluted rather than improved the best single model in this setting.

The TTA uncertainty is informative only for the convolutional networks (Table~\ref{tab:uncertainty}, Figure~\ref{fig:uncertainty}). For UNet, the nnU-Net-style network and VNet, cases with higher mean uncertainty inside the predicted tumour have lower Dice, with Spearman correlations of $-0.30$ to $-0.40$ and $p\leq 0.023$; for the three transformer-based models the correlation is between $-0.08$ and $-0.14$ and not significant.

\begin{heretable}
\centering
\caption{Spearman correlation between per-case TTA uncertainty inside the predicted WT and per-case mean Dice.}
\label{tab:uncertainty}
\small
\begin{tabular*}{\textwidth}{@{\extracolsep{\fill}}lcc@{}}
\toprule
Model & Spearman $\rho$ & $p$ \\
\midrule
UNet & $-0.3034$ & 0.0230 \\
nnU-Net & $-0.3040$ & 0.0228 \\
VNet & $-0.4018$ & 0.0021 \\
UNETR & $-0.0936$ & 0.4924 \\
Swin UNETR & $-0.0796$ & 0.5596 \\
TransBTS & $-0.1401$ & 0.3030 \\
\bottomrule
\end{tabular*}
\end{heretable}

\begin{herefigure}
\centering
\plot[0.7]{0.4}{segmentation/plots/09_uncertainty_vs_dice.png}
\caption{Per-case TTA uncertainty against per-case mean Dice for every model.}
\label{fig:uncertainty}
\end{herefigure}

\subsection{Computational Efficiency}
\label{sec:efficiency}

Table~\ref{tab:efficiency} compares cost. Parameter count does not predict accuracy: UNETR, with 137.60 million parameters, has the lowest Dice, whereas VNet, with 1.40 million, is third and reaches 97\% of the UNet Dice (0.7965 against 0.8199). VNet is also the cheapest model on every axis: 149.9 GFLOPs per volume, 0.53 training hours, 3.44~GB peak training memory, 0.98~GB at inference and 0.0375~s per volume without TTA, giving the best Dice per million parameters (0.5694) and per training hour (1.4939). UNet costs about three times as much per volume (0.1076~s) and 2.4 times the training time for its 0.0234 Dice advantage. Swin UNETR is the slowest model at 0.4893~s per volume and 5.29 training hours, despite having only 10.52 million parameters; with $2\times 2\times 2$ patches, its first stage attends over $64^3$ tokens. TTA multiplies latency by about eight for every model.

\begin{heretable}
\centering
\caption{Computational cost. GFLOPs are estimated per $128^3$ volume; latency is per volume at batch size 1 under half-precision autocasting, mean $\pm$ standard deviation; memory is peak GPU memory in GB for training and inference.}
\label{tab:efficiency}
\small
\setlength{\tabcolsep}{3pt}
\begin{tabular*}{\textwidth}{@{\extracolsep{\fill}}lcccccccc@{}}
\toprule
 & & & & \multicolumn{2}{c}{Memory (GB)} & \multicolumn{2}{c}{Latency (s)} & \\
\cmidrule(lr){5-6}\cmidrule(lr){7-8}
Model & Params (M) & GFLOPs & Train (h) & Train & Infer. & No TTA & TTA & Dice / M \\
\midrule
UNet & 5.60 & 907.0 & 1.28 & 5.25 & 1.54 & $0.1076\pm 0.0013$ & $0.8611\pm 0.0019$ & 0.1463 \\
nnU-Net & 16.55 & 956.3 & 1.35 & 5.22 & 1.65 & $0.1112\pm 0.0017$ & $0.8963\pm 0.0021$ & 0.0495 \\
VNet & 1.40 & 149.9 & 0.53 & 3.44 & 0.98 & $0.0375\pm 0.0003$ & $0.2973\pm 0.0007$ & 0.5694 \\
UNETR & 137.60 & 1761.7 & 2.49 & 9.30 & 3.80 & $0.2094\pm 0.0029$ & $1.6966\pm 0.0048$ & 0.0051 \\
Swin UNETR & 10.52 & 797.1 & 5.29 & 7.07 & 6.97 & $0.4893\pm 0.0034$ & $3.9512\pm 0.0044$ & 0.0709 \\
TransBTS & 35.81 & 996.2 & 1.56 & 7.46 & 1.99 & $0.1291\pm 0.0010$ & $1.0336\pm 0.0022$ & 0.0204 \\
\bottomrule
\end{tabular*}
\end{heretable}

Dice is plotted against latency in Figure~\ref{fig:latency}.

\begin{herefigure}
\centering
\plot[0.7]{0.4}{segmentation/plots/10_dice_vs_latency.png}
\caption{Test mean Dice against latency per volume.}
\label{fig:latency}
\end{herefigure}

\subsection{Discussion}
\label{sec:discussion}

Under one split, one training length and one optimisation budget, the three convolutional networks form a clearly better group than the three transformer-based models, and the gap is significant for every pair. Several factors are consistent with this outcome, although the experiments do not isolate them. Convolutions carry locality and translation equivariance by construction, which matters with 257 training volumes, whereas the transformer encoders here are trained from scratch without large-scale pretraining \citep{touvron2021}. UNETR tokenises the volume into $16^3$ patches, so its 512 tokens are coarse relative to small tumour sub-regions, which is consistent with its weak TC boundaries. The transformer-based models are also markedly more sensitive and less precise, so a different threshold or stronger post-processing might narrow the gap; the thresholds here were fixed at 0.5 for all models.

The nnU-Net-style network matches UNet but does not improve on it: its five stages and extra deep-supervision output cost three times the parameters for a median per-case difference of 0.0023 Dice. This network is not the self-configuring nnU-Net pipeline, so the result says nothing about nnU-Net's automatic configuration. VNet is the clear efficiency choice and is the model deployed in the accompanying web application, where its exported checkpoint (best validation Dice 0.7598, test mean Dice 0.7965) segments a four-sequence case in a single forward pass. Mean Dice should be read together with the per-region and per-grade results: every model is reliable for the whole tumour, while enhancing-tumour segmentation of low-grade gliomas is the weakest result of every architecture.

\section{Limitations and Future Work}
\label{sec:limitations}

Several limitations affect this study. All results come from a single dataset, the BraTS 2020 training cases, with no external multi-centre test set with public labels, so they may not generalise to other scanners and protocols. Every model was trained once with one seed, so seed-to-seed variation is not measured; the test set has 56 cases, of which only 12 are LGG. The transformer encoders received no large-scale pretraining. A single input resolution of $128^3$ was used for the whole resampled volume, which discards detail at native resolution. Hyperparameters followed one recipe per family rather than being tuned per model.

Future work should repeat the benchmark over several seeds and with cross-validation, add pretrained transformer encoders, train at higher resolution or with patch-based sampling, tune post-processing for each model, evaluate on external cohorts, and study LGG cases without enhancing tumour specifically, for example with grade-aware post-processing.

\section{Ethical Implications}
\label{sec:ethics}

Clinical deployment risks must be carefully considered. Over-reliance on AI predictions could lead to missed diagnoses if they replace rather than supplement radiologist judgment. False negatives could delay treatment, and false positives from over-segmentation might result in unnecessary interventions. The results by grade show that every model segments the enhancing tumour of low-grade gliomas far worse than that of high-grade gliomas, a bias that must be disclosed to users. The BraTS dataset lacks diversity in age, gender, and ethnicity, with data primarily originating from North American and European institutions, so models may not generalise to settings with different imaging protocols, and segmentation masks reflect individual annotator expertise. Responsible deployment requires human-in-the-loop review of all AI predictions, continuous monitoring of model performance on new data, uncertainty estimates with each prediction, regulatory compliance, and adherence to data privacy standards. The released code is intended for research and must not be used for autonomous diagnosis without clinical validation.

\section{Conclusion}
\label{sec:conclusion}

This study compared six 3D segmentation architectures on BraTS 2020 under one controlled protocol: identical preprocessing, a grade-stratified split, 50 epochs and the same number of optimiser steps for every model, and per-case evaluation at native resolution with confidence intervals and Holm-corrected significance tests. UNet achieved the highest mean Dice of 0.8199, the nnU-Net-style network 0.8192 and VNet 0.7965, ahead of Swin UNETR (0.7462), TransBTS (0.7291) and UNETR (0.7056), and every pairwise difference was significant. Whole-tumour segmentation was good for every model, but enhancing-tumour Dice on low-grade gliomas fell to 0.16--0.42. Parameter count did not predict performance: VNet, with 1.40 million parameters, reached 97\% of the UNet Dice at a third of its latency and less than half its training time. For the highest accuracy, UNet is recommended; for limited computational resources and deployment, VNet; transformer-based models should be reconsidered when pretraining, larger datasets or higher-resolution training are available.

\section*{Declaration of competing interest}

The authors declare that they have no known competing financial interests or personal relationships that could have appeared to influence the work reported in this paper.

\section*{Acknowledgements}

This research was conducted using Kaggle's GPU resources.

\section*{Data availability}

The BraTS 2020 training data are publicly available (Section~\ref{sec:dataset}). The notebook, the split of every case, the per-run histories, the per-case test metrics and the exported VNet checkpoint are available at the repository given in \ref{app:code}.

\appendix

\section{Code Repository}
\label{app:code}

The notebook that performs preprocessing, training, evaluation and figure generation, together with its outputs (the split of every case, per-run histories, per-case test metrics, plots and the exported VNet checkpoint), is available at \url{https://github.com/parthmital/NeuroScan-AI}, in the folder \texttt{Materials/Notebooks}.

\bibliographystyle{cas-model2-names}
\bibliography{refs}

\end{document}
